% Pista ana-24i-q1 · Anàlisi
% Procedència: PAU Catalunya 2024 · Convocatòria de juny (incidències) · Sèrie 5 · Qüestió 1
\documentclass[11pt,a4paper]{article}
\usepackage{pista}
\pistainfo{Catalunya 2024}{Convocatòria de juny (incidències)}{Sèrie 5}

\begin{document}

\section*{\textcolor{blauPista}{Qüestió 1}}

\noindent Considereu la funció $f(x) = -2 + 10\,(x-1)\ln x$, definida per a $x > 0$.

\subsection*{\textit{a)} \normalfont Comproveu que $f(x)$ té una arrel a l'interval $[1, 1{,}5]$ i busqueu un interval d'una dècima de longitud que també contingui aquesta mateixa arrel. \hfill\textnormal{[0,75 punts]}}

\pista{Aquí has d'aplicar el Teorema de Bolzano: si $f$ és contínua en l'interval i canvia de signe entre els extrems, llavors hi ha alguna arrel. Calcula $f(1)$ i $f(1{,}5)$ i comprova-ho. Per a trobar l'interval d'amplitud $0{,}1$ que conté l'arrel, prova valors intermedis ($1{,}1$, $1{,}2$, etc.) fins a localitzar un canvi de signe. O bé també pots fer, per exemple, comparar f(1) i f(1,25), i així vas fent la meitat cada vegada. Són dos mètodes alternatius, els dos són vàlids}

\subsection*{\textit{b)} \normalfont Sense calcular els punts crítics, justifiqueu que $f(x)$ és decreixent a l'interval $(0, 1)$ i creixent a $(1, +\infty)$. Quins màxims i mínims té aquesta funció? \hfill\textnormal{[1 punt]}}

\pista{La primera cosa: si ens diuen explícitament que no podem calcular punts crítics, vol dir que NO podem derivar i resoldre l'equació $f'(x)=0$. Hem de trobar altres maneres. Comencem amb algunes observacions. La constant $-2$ no afecta el creixement o decreixement d'una funció. I, si ho penses, la constant multiplicativa $10$ tampoc té cap influència en els intervals on la funció creix o decreix. Per tant, tot es redueix a estudiar el comportament del producte $(x-1)\ln x$. Observa els signes dels dos factors a cada interval: a $(0,1)$ són \emph{tots dos} negatius i a $(1, +\infty)$ són \emph{tots dos} positius. Analitza com varia cada factor dins de cada interval i dedueix-ne la monotonia del producte. Per als extrems, observa que $g(1) = 0$ i que $g(x) \geq 0$ a tot el domini: això et diu directament on hi ha el mínim (i si hi ha, o no, màxim).}

\subsection*{\textit{c)} \normalfont Calculeu $\displaystyle \lim_{x \to 0^{+}} f(x)$ i $\displaystyle \lim_{x \to +\infty} f(x)$, i feu un esbós de la gràfica d'aquesta funció. \hfill\textnormal{[0,75 punts]}}

\pista{Els dos límits es redueixen al comportament del producte $(x-1)\ln x$ (la constant $-2$ és irrellevant). Quan $x \to 0^{+}$, el factor $(x-1)$ tendeix a $-1$ i $\ln x$ tendeix a $-\infty$: d'aquí pots deduir el límit directament, sense eines avançades. Quan $x \to +\infty$, els dos factors tendeixen a $+\infty$. Per a l'esbós, combina aquests dos límits amb el creixement i el mínim de l'apartat b) i amb l'arrel localitzada a l'apartat a).}

\end{document}
